\section{Phân tích và Thảo luận (Discussion)}
\label{sec:discussion}

\subsection{Pipeline cũ đang bù đắp cho SciNCL}

Phát hiện quan trọng nhất của giai đoạn đầu là \textbf{toàn bộ kiến trúc đa tầng --- RRF, selective
reranking, QPP, CRC --- thực chất đang bù đắp cho một dense retriever yếu}. SciNCL chỉ đạt $\ndcg = 0.5640$
(MAP@10 = 0.5077), buộc hệ thống phải dùng BM25 ($\ndcg = 0.6523$) và RRF để đạt 0.6583, rồi
cross-encoder để đẩy lên 0.6939. Mỗi tầng thêm vào đều tăng chi phí và độ phức tạp. Khi thay SciNCL bằng
BGE-base, bức tranh hoàn toàn đảo ngược: BGE-base đơn thuần ($\ndcg = 0.7376$) đã vượt qua toàn bộ
pipeline cũ ($+0.1737$, $p<0.001$). BM25 và RRF trở nên không cần thiết, và cross-encoder \emph{luôn
làm giảm} chất lượng trên BGE pipeline ($-0.0374$, $p=0.015$) do không gian embedding khác biệt ---
MiniLM-L-6-v2 được huấn luyện trên MS MARCO web search, trong khi BGE được huấn luyện trên dữ liệu
đa dạng bao gồm văn bản khoa học qua MTEB-style contrastive learning.

\subsection{Specialization-Generalization Trade-off}

Fine-tune BGE-small (33M, 384d) trên SciFact tạo ra một specialist cực mạnh: $\ndcg = 0.7909$ --- vượt
BGE-base (109M, 768d) tới $+0.0533$, và vượt cả vstash reference cùng kích thước ($+0.0202$). Đây là
minh chứng rằng self-supervised fine-tuning với RRF disagreement signal có thể tạo ra mô hình nhúng
vượt trội trên miền chuyên biệt từ tập dữ liệu nhỏ.

Tuy nhiên, cái giá phải trả là transfer performance: trên ba bộ dữ liệu ngoài miền (NFCorpus, FiQA,
SciDocs), BGE-small specialist chỉ đạt transfer average $0.3011$, kém BGE-base generalist tới $-0.0239$.
Đây là một minh họa kinh điển của \emph{catastrophic forgetting} trong fine-tuning: mô hình nhỏ (33M)
mất khả năng tổng quát nhanh hơn mô hình lớn (109M) do dung lượng tham số thấp hơn.

\subsection{Từ Adaptive Coverage đến Batch Mode-Switch}

Câu hỏi then chốt: làm thế nào để \emph{giữ} specialist trên SciFact trong khi \emph{phục hồi} transfer
trên các miền khác? Các phương pháp per-query coverage (A3 fixed 5\%, adaptive uncertainty coverage)
đều không giải quyết được triệt để: chúng quá bảo thủ (coverage thấp) hoặc quá nhiễu (user feedback
mechanism quá mạnh), dẫn đến transfer vẫn kém BGE-base đáng kể ($-0.0152$).

\textbf{B5 giải quyết điểm nghẽn này bằng batch-level mode switch.} Thay vì quyết định rescue ở mức
từng query --- vốn dễ bị nhiễu --- B5 tính batch shift score $S$ từ trung bình z-score của năm đặc trưng
trên toàn batch, và quyết định \emph{toàn bộ batch} nên ở specialist-safe hay generalist-fallback. Cơ
chế này có ba ưu điểm:

\begin{enumerate}[leftmargin=2em,itemsep=2pt]
  \item \textbf{Giảm nhiễu:} Trung bình qua batch làm mượt các outlier query riêng lẻ, giúp quyết định ổn
  định hơn.
  \item \textbf{Frozen statistics:} Tất cả z-score được chuẩn hóa bằng $\mu_{\text{trainfit}}$ và
  $\sigma_{\text{trainfit}}$ từ SciFact --- không dùng label của target dataset, do đó không bị ảnh
  hưởng bởi train/test distribution shift.
  \item \textbf{Phân tách rõ ràng:} SciFact có batch shift score thấp ($S \approx 1.09$), trong khi
  NFCorpus ($6.03$), FiQA ($3.79$), và SciDocs ($3.58$) đều có $S$ cao --- controller phân biệt được
  source batch và transfer batch một cách đáng tin cậy.
\end{enumerate}

Kết quả: Frozen B5 gần như phục hồi hoàn toàn transfer average của BGE-base ($0.3249$ vs $0.3251$,
$\Delta = -0.0001$), trong khi SciFact vẫn giữ ở $0.8218$ ($+0.0030$ so với BGE-small). Đây là bước nhảy
lớn nhất trong toàn bộ hành trình phát triển ($+0.0151$ từ current adaptive SGAF).

\subsection{P3 Rank-Window Smoothing: Tín hiệu Specialist trong Top Window}

P3 trả lời câu hỏi: \emph{khi B5 đã chọn BGE-base làm ranking chính, liệu BGE-small còn cung cấp tín
hiệu hữu ích nào không?} Kết quả cho thấy có --- nhưng chỉ trong một cửa sổ nhỏ ở đầu ranking.

P3 chỉ kích hoạt trong generalist-fallback batches. Nó pha một prior BGE-small yếu ($\alpha = 0.10$) vào
top-20 của ranking BGE-base qua weighted RRF. Ý tưởng: BGE-small tuy yếu hơn trên transfer, nhưng vẫn
có thể xác định được một số document chất lượng mà BGE-base bỏ sót ở vị trí cao. Bằng cách blend nhẹ
hai ranking trong top-20, P3 tận dụng tín hiệu bổ trợ này mà không làm nhiễu phần đuôi.

P3 cải thiện transfer average từ $0.3249$ lên $0.3293$ ($+0.0043$), vượt BGE-base $+0.0042$. Cải thiện
này dương trên toàn bộ 6 dataset (4 nội bộ + 2 external), với significance trên NFCorpus ($p=0.02$) và
ArguAna ($p=0.02$). Quan trọng: P3 không bao giờ làm giảm kết quả so với B5 trên bất kỳ dataset nào
trong tập đánh giá hiện tại.

Leave-one-transfer-dataset-out (LOTO) validation chọn cùng bộ tham số $(w=20, \alpha=0.10)$ trong cả ba
fold, và mọi held-out delta đều dương --- giảm lo ngại về overfit trên grid exploration.

\subsection{External Validation và Hạn chế của Shift Detector}

Kết quả external validation trên TREC-COVID và ArguAna cho thấy hai mặt của B5/P3. Trên TREC-COVID
(biomedical, gần với SciFact), B5/P3 hoạt động tốt: P3 đạt $\ndcg = 0.6908$, cao hơn BGE-base $+0.0073$.
Tuy nhiên, trên ArguAna (argument retrieval, khác biệt lớn), B5 dưới BGE-base $0.0229$ --- shift
detector huấn luyện trên SciFact không phát hiện được argument retrieval batches là "khác biệt", dẫn đến
một số batch bị giữ ở specialist-safe mode trong khi đáng lẽ nên fallback. P3 bù đắp một phần ($+0.0011$,
$p=0.02$) nhưng không thể đóng hoàn toàn khoảng cách.

Bài học: SciFact-fitted shift detector không khái quát tốt sang argument retrieval. Đây là hướng cải
thiện trong tương lai: domain-adaptive shift detection hoặc multi-source calibration.

\subsection{Luận điểm ``Cheap First''}

Toàn bộ cải thiện từ B5/P3 đến từ các thao tác post-retrieval CPU thuần, với chi phí bổ sung dưới
$0.01$\,ms/query. So với cross-encoder ($27.5$\,ms/query, GPU), B5/P3 rẻ hơn $\sim 2700\times$ trên mỗi
query. So với toàn bộ pipeline BGE retrieval ($2.6$\,ms/query), B5/P3 chỉ thêm chưa đến $0.4\%$ latency.

Bảng~\ref{tab:cost-analysis} cho thấy B5/P3 đạt transfer average cao nhất ($0.3293$) với latency thấp thứ
hai ($2.6$\,ms/q, chỉ sau BGE-small/Base đơn thuần ở $2.3$\,ms/q). Luận điểm ``cheap first'' của chúng
tôi là: \textbf{trước khi thêm GPU inference (cross-encoder, LLM), hãy khai thác tối đa tín hiệu đã có
sẵn từ retrieval bằng các thao tác post-retrieval CPU}. B5/P3 chứng minh rằng cách tiếp cận này có thể
phục hồi transfer và thậm chí vượt generalist baseline --- với chi phí gần như bằng không.

Oracle headroom (transfer average của Oracle Router = $0.3855$) cho thấy vẫn còn dư địa cải thiện đáng
kể, nhưng để đạt được mức đó cần các phương pháp đắt tiền hơn (cross-encoder, LLM reranker). Điều này
củng cố luận điểm: B5/P3 là giải pháp ``cheap repair'' hiệu quả nhất trước khi chuyển sang các tầng
đắt tiền.

\subsection{Ý nghĩa thống kê và Độ tin cậy}

Chúng tôi tiến hành paired bootstrap (10\,000 resamples) trên nDCG@10 cho tất cả so sánh SGAF. Kết quả:

\begin{itemize}[leftmargin=1.6em,itemsep=2pt]
  \item \textbf{Frozen B5 vs BGE-small:} Có ý nghĩa trên NFCorpus ($+0.0187$, $p=0.01$), FiQA ($+0.0274$,
  $p=0.001$), và SciDocs ($+0.0253$, $p<0.001$).
  \item \textbf{Frozen P3 vs Frozen B5:} Significant trên NFCorpus ($+0.0052$, $p=0.02$) và ArguAna
  ($+0.0011$, $p=0.02$); directionally positive nhưng chưa đạt significance trên FiQA ($p=0.10$) và
  SciDocs ($p=0.07$).
  \item \textbf{Frozen P3 vs BGE-base:} Dương trên toàn bộ transfer dataset nhưng chưa significant
  trên bất kỳ dataset riêng lẻ nào --- đúng với kỳ vọng: P3 cải thiện nhỏ nhưng nhất quán, cần nhiều
  query hơn để đạt significance ở mức dataset.
\end{itemize}

\subsection{Train/Test Distribution Shift}

Một phát hiện phụ quan trọng: SciFact có sự khác biệt phân phối rất lớn giữa train và test. Train oracle
labels thiên về Dense (477/809), trong khi test thiên về BM25 (226/300). Điều này khiến mọi phương pháp
học có giám sát trên SciFact trainfit --- learned QPP, LLM query routing, BM25 tuning --- đều thất bại
khi chuyển từ train sang test. Các phương pháp không giám sát (BGE embeddings, Adaptive RRF dùng IDF
corpus, B5/P3 với z-score frozen từ trainfit) miễn nhiễm với vấn đề này \emph{miễn là} thống kê dùng để
chuẩn hóa không bị ảnh hưởng bởi label shift.

B5/P3 né được vấn đề này bằng cách chỉ dùng $\mu_{\text{trainfit}}$ và $\sigma_{\text{trainfit}}$ của
các đặc trưng query (không phải label), và không huấn luyện tham số mới trên trainfit. Logistic
regression rescue ranker được huấn luyện trên trainfit nhưng chỉ dùng trong specialist-safe mode (batch
giống SciFact), nơi distribution shift nhỏ nhất.
